\subsection{Completed objects: abstract maps and natural families}

The manuscript uses two descriptions on scans 176 and 178: the
automorphisms of the completed algebra at a fixed ring, and the
automorphisms of its functor in the coefficient ring. These need
not mean the same thing. The following statements resolve this
distinction, without assuming automatic continuity.

\paragraph{B-15. Refuted: abstract automatic continuity and
classification by one-variable data.}
Over $k=\mathbf F_2$ there is an abstract simplicial differential
$k$-algebra automorphism $\Phi$ of $\widehat A_k$ which fixes every
polynomial form but is discontinuous. It has infinite order.
Moreover $\Phi^2$ is the identity in simplicial dimension one but
is not the identity in dimension two. Thus even allowing arbitrary
one-variable derivations in place of convergent $B_i$ does not
classify all abstract automorphisms.

\emph{Proof.}
Write $y_m(V)=V^{(2^m)}$. For each integer $N\geq0$, take the
continuous B-automorphism $\phi_N$ of B-03 with $A=X$, $c_N=1$,
and $c_m=0$ for $m\ne N$. Its one-variable derivation is
\[
 \delta_N(V^{(a)})=
 \begin{cases}
 V^{(a-2^N)},&\text{the $N$th binary digit of $a$ is one},\\
 0,&\text{otherwise}.
 \end{cases}
 \tag{B-A1}
\]
In a simplex with cumulative coordinates $V_1,\ldots,V_n$, let
$\delta_{N,i}$ act in $V_i$ and fix the other coordinates and $T$.
The exact finite expression is
\[
 \phi_N(f)=\sum_{I\subseteq\{1,\ldots,n\}}
        \left(\prod_{i\in I}\delta_{N,i}\right)(f)\,dV_I,
 \qquad \phi_N(dV_i)=dV_i.
 \tag{B-A2}
\]
The product of operators is unambiguous: these particular continuous
coordinate derivations commute. Formula (B-A2) has $2^n$ terms and
is the extension of the polynomial formula in B-01.

Choose a nonprincipal ultrafilter $\mathcal U$ on the nonnegative
integers. Each $\widehat A_{k,n}$ is the product of its finite
weight spaces, hence a compact Hausdorff space. For every element
$a$ define
\[
 \Phi_n(a)=\lim_{N\to\mathcal U}\phi_{N,n}(a).
 \tag{B-A3}
\]
Equivalently each coefficient is the unique value in $\mathbf F_2$
taken on a set belonging to $\mathcal U$. The same ultrafilter is
used for every element and every simplex. Addition, multiplication,
the differential, and each simplicial operator are continuous in
the weight topology: every output coefficient uses finitely many
input coefficients. Taking the limit in each of their defining
identities proves that $\Phi$ is an abstract unital simplicial
differential-algebra endomorphism.

Every $\phi_N$ is the identity on the associated graded for exterior
degree. The exterior filtration is closed. Consequently
$Q_n=\Phi_n-1$ strictly raises exterior degree and $Q_n^{n+1}=0$.
The finite sum $\sum_{j=0}^n(-Q_n)^j$ is its inverse as an additive
map. The inverse of a bijective differential-ring homomorphism is
again such a homomorphism, and the inverse maps commute with the
simplicial operators. Thus $\Phi$ is an automorphism.

Every polynomial form is fixed by $\phi_N$ for all sufficiently
large $N$, and therefore by $\Phi$. On the other hand, the series
\[
 s(V)=\sum_{m\geq0}y_m(V)
\]
satisfies $\phi_N(s)=s+dV$ for every $N$, so
$\Phi(s)=s+dV$. For the tails
$s_r=\sum_{m\geq r}y_m$, one has $s_r\to0$ but
$\Phi(s_r)=s_r+dV\not\to0$. This proves discontinuity without any
appeal to an uncomputed coefficient.

There is also an explicit witness for the higher-dimensional data.
Put $F_0=1$ and, in dimension $n$,
\[
 F_n(V_1,\ldots,V_n)=
   \sum_{0\leq m_1<\cdots<m_n}
          y_{m_1}(V_1)\cdots y_{m_n}(V_n).
 \tag{B-A4}
\]
This series converges in total weight. Applying (B-A2), no term can
have two removed indices equal to $N$, since the indices are
strictly increasing. If the removed index is not the last one,
the remaining later index is greater than $N$, so that term tends
to zero in each fixed weight. The terms with the last index removed
converge to $F_{n-1}dV_n$. Hence
\[
 Q(F_n)=F_{n-1}dV_n,\qquad
 Q^r(F_n)=F_{n-r}dV_{n-r+1}\cdots dV_n\quad(0\leq r\leq n).
 \tag{B-A5}
\]
In particular $\Phi^2(F_2)=F_2+dV_1dV_2$. In dimension one,
$Q^2=0$, so $\Phi_1^2=1$. More generally, in characteristic two,
$\Phi^{2^r}=1+Q^{2^r}$; its action on $F_{2^r}$ is not the identity.
Each fixed simplex has a power of two as an exponent for this
unipotent map, so these nonvanishings imply that $\Phi$ has infinite
order. This construction uses the ordinary ultrafilter principle;
it is an existence proof with explicit witnesses, not a finite
computational certificate.\hfill$\square$

\paragraph{B-16. Proved: continuity in degree zero and the abstract determinant criterion.}
Let $\Omega_{\rm abs}(k)$ mean automorphisms with no topology
assumed. Every underlying simplicial differential $k$-algebra
endomorphism of $\widehat A_k$ preserves the exterior filtration.
Its exterior-degree-zero part is the continuous substitution of
B-03 for a uniquely determined
$A\in(X)_{\rm PD}^{\wedge}\subset\widehat S\{\{X\}\}$ and
$u(T)=A(T,T)$. Such an endomorphism is invertible if and only if
$\Delta=a_{10}$ is a unit. In particular there is a split exact
sequence of groups
\[
 1\longrightarrow W(k)\longrightarrow\Omega_{\rm abs}(k)
       \longrightarrow P(k)\longrightarrow1,
 \tag{B-A6}
\]
where $W(k)$ consists of the abstract maps inducing the identity on
exterior associated graded. Its continuous subgroup is $E(k)$ of
B-04; B-15 shows $W(k)$ is strictly larger in characteristic two.

\emph{Proof.}
Write $u=\Phi_0$ and $\Phi_1(X)=A+B\,dX$. The zero endpoint gives
$A(0)=0$, and the differential gives $\Phi_1(dX)=A' dX$.
By simplicial naturality, $\Phi_n(dU_i)=A'(U_i)dU_i$. Every form is
a finite sum of coefficient series times exterior monomials.
Multiplicativity therefore makes $\Phi$ preserve the exterior
filtration, even though the coefficient series need not be limits
preserved by $\Phi$.

Let $\psi$ be its degree-zero part in dimension one. For every
series $f=\sum_{a\geq0}f_a X^{(a)}$, differential compatibility and
the zero endpoint give
\[
 \partial_X\psi(f)=A'\psi(\partial_X f),\qquad
 \psi(f)|_{X=0}=u(f_0).
 \tag{B-A7}
\]
These equations determine every coefficient of $\psi(f)$ by
induction on its $X$-degree, simultaneously for all $f$; the
coefficient of degree $r$ on the right uses only coefficients of
$\psi(\partial_X f)$ of degrees at most $r$. The resulting formula is
\[
 \psi(f)=\sum_{a\geq0}u(f_a)A^{(a)}.
 \tag{B-A8}
\]
This comparison series is well defined in the $X$-adic divided-power
topology because $A^{(a)}$ has $X$-order at least $a$. It satisfies
(B-A7), so the coefficient induction proves equality without
passing an infinite sum through the unknown map.

Now take $f(X)=\sum a_rX^{(r)}$ with $a_r\in k$. The other endpoint,
$X=T$, yields
\[
 u\left(\sum a_rT^{(r)}\right)
     =\sum a_r\bigl(A(T,T)\bigr)^{(r)}.
 \tag{B-A9}
\]
Thus $u$ itself is a continuous divided-power substitution. In
several variables use
$\partial_i\psi(f)=A'(U_i)\psi(\partial_i f)$ and the vertex at
which all $U_i$ vanish. Induction on total $U$-degree gives precisely
the simultaneous substitution $U_i\mapsto A(T,U_i)$ over $u$.
It is consequently continuous in total weight.

If $\Delta$ is a unit, that degree-zero substitution is invertible
by B-03. Composing with its inverse leaves $1+Q$, with $Q$ raising
exterior degree, and the same finite inverse as in B-15 applies.
Conversely an inverse has an inverse degree-zero substitution;
their linear coefficients multiply to one. The substitution with
all $B_i=0$ supplies the section in (B-A6).\hfill$\square$

\paragraph{B-17. Proved: natural families are precisely continuous automorphisms.}
For a fixed ring $k$, consider natural families, indexed by
commutative $k$-algebras $B$, of simplicial differential
$B$-algebra automorphisms of $\widehat A_B$. Evaluation at $B=k$
identifies their group with $\Omega_{\rm cont}(k)$ of B-03--B-04.
This is the precise automorphism functor in the coefficient ring;
abstract automorphisms at one isolated ring must not be substituted
for it.

\emph{Proof.}
In a fixed simplicial dimension the underlying set functor is
$B\mapsto B^{\mathcal M}$, where $\mathcal M$ is the set of
divided-power exterior monomials. This functor is represented by
the polynomial algebra $k[z_\mu:\mu\in\mathcal M]$. Consequently
every output coefficient of a natural transformation is a
polynomial in finitely many input coefficients: evaluate it on the
universal point $(z_\mu)$, and use naturality for its evaluation in
every $B$.

The maps are $B$-linear. Evaluating scalar linearity over the
polynomial ring with an additional indeterminate $t$ shows that
each such polynomial $P$ satisfies $P(tz)=tP(z)$. Independence of
polynomial monomials over any base ring forces $P$ to be a finite
linear combination of the $z_\mu$. For every finite output-weight
range there are only finitely many output coordinates; the union
of their finite input supports is finite. This is exactly
continuity in the total-weight topology.

Conversely a continuous $k$-linear map has this finite-support
matrix description. Extend its matrix coefficients to $B$.
The unit, product, differential, simplicial, and inverse identities
extend as well: the linear identities are checked on monomials,
and the bilinear product identity on pairs of monomials. Each
output identity contains only finitely many input coefficients.
This constructs the unique natural family extending a continuous
automorphism and proves the claimed group identification.
\hfill$\square$

The exact manuscript locators are
\texttt{MPL\_136\_142-188\_transcription.tex}, lines 1160--1194
(completion), 1269--1275 (the fixed-ring automorphism notation),
and 1366--1373 (the coefficient-ring functor description).
Original scans 174, 176, 178, and 180 were fetched as a selected
PDF and inspected. The definition has no explicit continuity
adjective. The functorial description, interpreted as B-17, gives
the formal meaning for which B-03 is exact; the isolated abstract
interpretation has the extra automorphisms of B-15. In either
interpretation the cohomology action is the same, by B-06.
